\begin{helper}[Linearity of Expectation for Product Model]
Fix $n \in \mathbb{N}$, a subset $I \subseteq \mathrm{Fin}(n)$, and parameters $\varepsilon \in \mathbb{R}$.
Fix $m_{sub} \in \mathbb{N}$, a probability $p_{sub} \in \mathbb{R}$ with $0 \le p_{sub} \le 1$, and an embedding $\pi : \mathrm{Fin}(n) \hookrightarrow \mathrm{Fin}(m_{sub}) \times \mathrm{Fin}(m_{sub})$.
Let $C = \noderef{RealChooseTwo}(|I|)$.
Suppose \(B\ge0\).  Suppose that for every pair of vertices \(u,w\in I\)
such that \(\pi(u)_1\ne\pi(w)_1\) and \(\pi(u)_2\ne\pi(w)_2\), the expected
value under the product model of the indicator of the following event is
bounded above by \(B\): either one of the first \(m_{sub}\) coordinates has
first component containing both \(\pi(u)_1,\pi(w)_1\), or one of the last
\(m_{sub}\) coordinates has second component containing both
\(\pi(u)_2,\pi(w)_2\).
Then the product expectation satisfies
\[
\sum_v \left(\prod_i q(v_i)\right) Z_{\mathrm{prod}}(v) \le C B,
\]
where $Z_{\mathrm{prod}}$ is given by $\noderef{FixedSetProductModelVar}$ and $q = \noderef{FixedSetProductModelMassFn}(p_{sub}, P_R, P_B)$.
\end{helper}
\begin{proof}
Let
\[
W(v)=\prod_{i\in\mathrm{Fin}(2m_{sub})}q(v_i).
\]
By \(\noderef{FixedSetProductModelMass}\), all one-coordinate masses are
nonnegative; hence \(W(v)\ge0\) for every assignment \(v\).

Fix an assignment \(v\) and unfold
\(\noderef{FixedSetProductModelVar}\).  The variable is the cardinality of a
finite set of pairs \(e=(u,w)\).  Each such pair is rank-oriented by the
final-pair construction used in this definition, both endpoints lie in \(I\)
because the relevant \(X\)-set is a filtered subset of \(I\), and the final
filter in \(\noderef{FixedSetProductModelVar}\) gives
\(\pi(u)_1\ne\pi(w)_1\) and \(\pi(u)_2\ne\pi(w)_2\).

Moreover, such a counted pair comes from some external base vertex.  If that
base vertex is red, then the red graph in
\(\noderef{FixedSetProductModelVar}\) is built from the first \(m_{sub}\)
coordinates, so for the corresponding coordinate \(x<m_{sub}\) the first
component \((v_x)_1\) contains both red projections \(\pi(u)_1,\pi(w)_1\).
If the base vertex is blue, then the blue graph is built from the last
\(m_{sub}\) coordinates, so for the corresponding coordinate
\(m_{sub}\le x<2m_{sub}\) the second component \((v_x)_2\) contains both
blue projections \(\pi(u)_2,\pi(w)_2\).  Thus every counted pair satisfies the
pair event appearing in the hypothesis.

Let \(E_{u,w}(v)\) be the indicator of that pair event.  The preceding
paragraph gives the pointwise counting inequality
\[
Z_{\mathrm{prod}}(v)
\le
\sum_{\{u,w\}\subseteq I} E_{u,w}(v),
\]
where unordered pairs are represented once by the rank orientation.  Multiplying
by the nonnegative weight \(W(v)\) preserves the inequality.  Summing over all
assignments and using finite distributivity gives
\[
\sum_v W(v)Z_{\mathrm{prod}}(v)
\le
\sum_{\{u,w\}\subseteq I}\sum_v W(v)E_{u,w}(v).
\]
For pairs whose red and blue projections are distinct, the inner sum is at
most \(B\) by hypothesis.  Pairs without distinct projections do not appear in
the pointwise upper bound for \(Z_{\mathrm{prod}}\), so their contribution can
be taken as \(0\), which is at most \(B\) because \(B\ge0\).  Therefore every
one of the \(\binom{|I|}{2}=\noderef{RealChooseTwo}(|I|)=C\) rank-oriented
unordered pairs contributes at most \(B\).  Hence
\[
\sum_v W(v)Z_{\mathrm{prod}}(v)\le C B.
\]
\end{proof}
